\subsection{Sobolev regularity at divisorial poles}
\label{subsec:divisorial-sobolev}\index{Sobolev regularity!divisorial poles}

At a divisorial pole, the rate of a convex transform determines which derivatives are integrable. The gradient and the classical Laplacian have different thresholds. The distributional Laplacian also retains a divisor measure unless the slope of the transform tends to zero.

Let $U\subset\mathbb C^n$ be a sufficiently small neighborhood of $0$, and let $f$ be a nonzero holomorphic function with $f(0)=0$. Write
\[
 f=u\prod_{j=1}^N h_j^{m_j},\qquad
 m_0=\min_jm_j,
\]
where $u$ is a holomorphic unit and the $h_j$ are distinct irreducible germs represented on $U$. Fix a convex nondecreasing $\chi\in C^2(({-\infty},A])$ and shrink $U$ so that $\log|f|<A-1$. Set $v=\chi(\log|f|)$ off $Z=\{f=0\}$, with its canonical plurisubharmonic extension across $Z$. All gradients, Laplacians, volumes and Sobolev spaces in this subsection are real Euclidean ones. Write $|\partial f|^2=\sum_j|\partial f/\partial z_j|^2$. Then $|\nabla|f||=|\partial f|$ off $Z$.

\Needspace{10\baselineskip}
\begin{lemma}\label{lem:holomorphic-small-level-energy}
After shrinking $U$, there is $C>0$ such that
\[
 |\partial f|\leq C|f|^{1-1/m_0}\quad\text{on }U\setminus Z.
\]
For each $K\Subset U$, there are $C_K,r_K>0$ with
\[
 \int_{K\cap\{|f|=r\}}|\partial f|\,\mathrm dS\leq C_Kr
 \qquad(0<r<r_K),
\]
where $\mathrm dS$ is real hypersurface measure. In particular, for nonnegative measurable $\Psi$,
\[
 \int_{K\cap\{0<|f|<r_K\}}\Psi(|f|)|\partial f|^2\,\mathrm dV
 \leq C_K\int_0^{r_K}\Psi(r)r\,\mathrm dr.
\]
\end{lemma}
\begin{proof}
Shrink $U$ so that $|h_j|\leq1$ and the unit, its inverse and all first derivatives are bounded. For each $j$,
$|f|^{1/m_0}\leq C|h_j|^{m_j/m_0}\leq C|h_j|$.
Logarithmic differentiation of the factorization therefore gives
\[
 \frac{|\partial f|}{|f|}
 \leq C\left(1+\sum_j|h_j|^{-1}\right)
 \leq C|f|^{-1/m_0},
\]
which proves the first bound.

For small $\zeta\in\mathbb C$, let $\mu_\zeta$ be the Euclidean volume measure of the fiber divisor $\{f=\zeta\}$, with multiplicity. The local compactness recalled in \cref{lem:qpsh-hartogs} and its preceding discussion gives $\log|f-\zeta|\to\log|f|$ in $L^1_{\mathrm{loc}}$ as $\zeta\to0$: every subsequential limit is identified off $Z$, where the convergence is locally uniform, and the same set excludes convergence to $-\infty$. The Poincar\'e--Lelong formula implies weak convergence of the fiber currents, hence of their Euclidean traces. Pairing with a fixed nonnegative cutoff equal to one near $K$ shows $\sup_{|\zeta|<r_K}\mu_\zeta(K)<\infty$.

The critical locus of $f$ has finitely many local components, and $f$ is constant on each. After shrinking $U$, every remaining critical component passes through $0$ and has value zero. Thus the small nonzero fibers are smooth. On $\{|f|=r\}$ the coarea formula for the map $f$ to the circle gives
\[
 \int_{K\cap\{|f|=r\}}|\partial f|\,\mathrm dS
 =\int_{|\zeta|=r}\mu_\zeta(K)\,\mathrm dS(\zeta)
 \leq C_Kr.
\]
Scalar coarea for $|f|$, using $|\nabla|f||=|\partial f|$, proves the last inequality. The same formulas hold in dimension one, with fiber volume interpreted as counting measure.
\end{proof}

\Needspace{12\baselineskip}
\begin{proposition}\label{prop:divisorial-first-sobolev}\index{Sobolev regularity!sharp first-order criterion}
For each $2\leq p<\infty$,
\[
 v\in W^{1,p}_{\mathrm{loc}}(U)
 \quad\Longleftrightarrow\quad
 \int_{-\infty}^A|\chi'(s)|^p
          \mathrm e^{(2-p)s/m_0}\,\mathrm ds<\infty.
\]
When this holds, its weak gradient is the zero extension of
$\chi'(\log|f|)\nabla\log|f|$ from $U\setminus Z$.
\end{proposition}
\begin{proof}
At a smooth point of a component of multiplicity $m_0$, outside the other components, a holomorphic root of the local unit changes the normal coordinate so that $f=w_1^{m_0}$. Bounded coordinate derivatives and Jacobians preserve local first-order Sobolev integrability. With $r=|w_1|$ and $s=m_0\log r$, polar integration of the classical gradient gives, up to fixed positive factors,
\[
 \int_0^{r_0}|\chi'(m_0\log r)|^p r^{1-p}\,\mathrm dr
 =\frac1{m_0}\int_{-\infty}^{m_0\log r_0}
       |\chi'(s)|^p\mathrm e^{(2-p)s/m_0}\,\mathrm ds.
\]
Such smooth points occur in every neighborhood of $0$. This proves necessity; changing the finite upper endpoint does not affect convergence.

For sufficiency put
\[
 s_\varepsilon=\tfrac12\log(|f|^2+\varepsilon^2),
 \qquad v_\varepsilon=\chi(s_\varepsilon).
\]
For small $\varepsilon>0$ these are $C^2$ plurisubharmonic functions decreasing to $v$. Convexity gives $|\chi(s)|\leq C(1+|s|)$ for $s\leq A$; local exponential integrability of $\log|f|$, recalled in \cite{SGPT}, therefore gives convergence in every finite local $L^p$ space. By the chain rule,
\[
 |\nabla v_\varepsilon|
 =\chi'(s_\varepsilon)
       \frac{|f|}{|f|^2+\varepsilon^2}|\partial f|.
\]
\Needspace{7\baselineskip}
Use \cref{lem:holomorphic-small-level-energy} in scalar coarea, leaving two powers of $|\partial f|$ for the level estimate and bounding the other $p-2$ powers by its first inequality. The resulting local gradient integral is bounded by
\[
 C\int_0^{r_K}
  |\chi'(\tfrac12\log(r^2+\varepsilon^2))|^p
  \frac{r^{p+(p-2)(1-1/m_0)+1}}{(r^2+\varepsilon^2)^p}
       \,\mathrm dr.
\]
With $s=\tfrac12\log(r^2+\varepsilon^2)$, use $\mathrm ds=r\,\mathrm dr/(r^2+\varepsilon^2)$ and $r\leq(r^2+\varepsilon^2)^{1/2}$ to bound this by
\[
 C\int_{-\infty}^A|\chi'(s)|^p
         \mathrm e^{(2-p)s/m_0}\,\mathrm ds.
\]
The part where $|f|\geq r_K$ is uniformly bounded as well. Weak Sobolev compactness and the local $L^p$ convergence identify the limit as $v\in W^{1,p}_{\mathrm{loc}}$. Its weak gradient agrees with the classical one off $Z$, and $Z$ has Lebesgue measure zero, proving the final assertion.
\end{proof}

\Needspace{13\baselineskip}
\begin{proposition}\label{prop:divisorial-laplacian-residual}\index{Laplacian!residual divisor measure}
Let $G$ be the zero extension of
$\chi''(\log|f|)|\nabla\log|f||^2$ from $U\setminus Z$, and put $\ell=\lim_{s\to-\infty}\chi'(s)$. For each $1\leq p<\infty$,
\[
 G\in L^p_{\mathrm{loc}}(U)
 \quad\Longleftrightarrow\quad
 \int_{-\infty}^A|\chi''(s)|^p
        \mathrm e^{2(1-p)s/m_0}\,\mathrm ds<\infty.
\]
As distributions,
\[
 \Delta v=G+\ell\,\Delta\log|f|.
\]
The second term is the positive divisor measure with coefficient $\ell$. In particular, $\Delta v$ is represented by a local $L^p$ function exactly when $\ell=0$ and the displayed integral is finite.
\end{proposition}
\begin{proof}
The limit $\ell\geq0$ exists because $\chi'$ is nonnegative and nondecreasing. Near a smooth divisor point, the coordinate reduction to $f=w_1^{m_0}$ makes the positive complex Hessian a single normal factor. Its Euclidean traces in the two coordinate systems are uniformly comparable, so polar integration gives necessity. In general,
$G=\chi''(\log|f|)|\partial f|^2/|f|^2$ off $Z$.
Scalar coarea and \cref{lem:holomorphic-small-level-energy}, now bounding the remaining $2p-2$ powers of $|\partial f|$, give
\[
 \int_{K\cap\{0<|f|<r_K\}}G^p\,\mathrm dV
 \leq C\int_0^{r_K}|\chi''(\log r)|^p
             r^{-1-2(p-1)/m_0}\,\mathrm dr.
\]
The substitution $s=\log r$ proves sufficiency. For $p=1$ the integral is always finite, since $\int_{-\infty}^A\chi''=\chi'(A)-\ell$.

\Needspace{7\baselineskip}
For the distribution identity, use $s_\varepsilon,v_\varepsilon$ from the preceding proof. Direct differentiation gives
\[
 \begin{aligned}
 \Delta v_\varepsilon
 &=\chi''(s_\varepsilon)
       \frac{|f|^2|\partial f|^2}{(|f|^2+\varepsilon^2)^2}
       +\chi'(s_\varepsilon)\Delta s_\varepsilon,\\
 \Delta s_\varepsilon
 &=\frac{2\varepsilon^2|\partial f|^2}
         {(|f|^2+\varepsilon^2)^2}.
 \end{aligned}
\]
\Needspace{5\baselineskip}
Off $Z$, the first term tends locally uniformly to $G$. Its integral over $K\cap\{|f|<r_0\}$ is at most
\[
 C\int_{-\infty}^{\frac12\log(r_0^2+\varepsilon^2)}
                  \chi''(s)\,\mathrm ds,
\]
by the level-energy estimate and the same substitution. First let $\varepsilon\to0$ and then $r_0\to0$. The bound tends to zero, so this term converges locally in $L^1$ to $G$.

The positive measures $\Delta s_\varepsilon$ converge weakly to $\Delta\log|f|$, with locally bounded masses. The Poincar\'e--Lelong formula identifies this limit as a measure supported on $Z$. On $\{|f|<r_0\}$,
\[
 0\leq\chi'(s_\varepsilon)-\ell
 \leq\chi'(\tfrac12\log(r_0^2+\varepsilon^2))-\ell;
\]
on compact subsets where $|f|\geq r_0$, $\Delta s_\varepsilon\to0$ uniformly. The same successive limits show that
$(\chi'(s_\varepsilon)-\ell)\Delta s_\varepsilon\to0$ weakly. Passing to distributions gives the identity. Since the nonzero divisor measure is singular with respect to Lebesgue measure, uniqueness of the Lebesgue decomposition proves the final criterion.
\end{proof}

\Needspace{12\baselineskip}
\begin{example}\label{ex:divisorial-sobolev-endpoint}\index{Sobolev regularity!endpoint failure}
Shrink $U$ so that $\log|f|<-1$. For $0<a<1$, the plurisubharmonic function
\[
 v_a=-(-\log|f|)^a
\]
has $\Delta v_a\in L^1_{\mathrm{loc}}$, while
\[
 v_a\in W^{1,2}_{\mathrm{loc}}
       \quad\Longleftrightarrow\quad a<\tfrac12.
\]
It belongs to no $W^{1,p}_{\mathrm{loc}}$ with $p>2$ and to no $W^{2,1}_{\mathrm{loc}}$.
\end{example}
\begin{proof}
Here $\chi'(s)=a(-s)^{a-1}$ and $\chi''(s)=a(1-a)(-s)^{a-2}$, with $\ell=0$. The preceding two propositions give the Laplacian assertion and the first-order thresholds. At $a=1/2$ the first-order integral diverges logarithmically; for $p>2$ its exponential factor forces divergence.

For second derivatives, take a smooth divisor point with normal coordinate $w_1$ and write $F(r)=\chi(m\log r)$, $r=|w_1|$. The real tangential Hessian in the unit direction along the circle has value $F'(r)/r=m\chi'(m\log r)/r^2$. Thus its absolute integral in the normal disc is bounded below, up to a fixed positive factor, by
\[
 \int_0^{r_0}\chi'(m\log r)\frac{\mathrm dr}{r}
 =\frac1m\int_{-\infty}^{m\log r_0}\chi'(s)\,\mathrm ds
 =+\infty.
\]
The classical Hessian off the divisor therefore fails local $L^1$ integrability, which is necessary for $W^{2,1}_{\mathrm{loc}}$. Smooth coordinate changes preserve that Sobolev class, completing the argument.
\end{proof}
